\documentclass[a4paper,10pt]{article}
\usepackage[margin=1.5cm]{geometry}

\begin{document}
\section*{Python en  campos de la tecnología}
Al ser un un lenguaje de proposito general, \textbf{Python}
se ha convertido en una buena alternativa en campos como
    \begin{itemize}
        \item Análisis de datos: Herramientas como \textbf{Pandas} y \textbf{Numpy} 
        hace que el Análisis de grandes conjuntos de datos sea ménos tedioso y accesible.
        \item Machine Learning: Tecnologias como \textbf{TensorFlow, Sckit-learn y Pytorch},
        permiten a los desarrolladores trabajar con modelos de IA de forma más accesible y directa.
        \item Desarrollo Web: Frameworks como \textbf{Django, FastAPI y FLASK} permiten a los dessarrolladores
        implementar sistemas \textbf{backend} seguros de forma sencilla.
        \item Ciberseguridad: Python es muy utilizado en el campo para detectar vulnerbilidades, malware y virus. También
        permite crear sistemas de análisis automatizados.
        \item Sistemas embebidos: Python corre bien en microcontroladores como \textbf{Raspberry Pi} y
        en placas compatibles com \textbf{MicroPython}. Esto permite crear dispositivos inteligentes.
        \item Web Scrapping: Librerias como \textbf{Selenium y BeautifulSoup} vuelven sencillo interactuar con 
        sitios web, permitiendo todo tipo de automatización a través de la web (tareas, Data Scrapping, despliegues web, etc.).
    \end{itemize}

Dicho esto, Python es un lenguje de programación muy versatil, que ofrece diferentes alternativas para las variadas ramas de la
computación.



\section*{Tecnologías compatibles con Python}
\begin{center}
    \begin{tabular}{| c | c | c |}
        \hline
        Tecnología & Framework/Librería & Descripción \\
        \hline\hline
    \end{tabular}
\end{center}


\end{document}